\documentclass[11pt]{article}

\usepackage[T1]{fontenc}
\usepackage[utf8]{inputenc}
\usepackage{lmodern}
\usepackage[a4paper,margin=27mm]{geometry}
\usepackage{microtype}
\usepackage{amsmath,amssymb,mathtools}
\usepackage{enumitem}
\usepackage{booktabs,tabularx,longtable}
\usepackage[colorlinks=true,linkcolor=blue,citecolor=blue,urlcolor=blue]{hyperref}

\title{A Projection-First Reframing of String Theory:\\
\large Conditional Encodings, Worldsheet Gates, and the q79 Boundary}
\author{Peter Nero}
\date{July 2026\\Version 2}

\begin{document}
\maketitle

\begin{abstract}
String theory supplies a highly constrained language for extended
excitations, quantum gravity, gauge sectors, duality, and compactification.
A projection-first viewpoint can organize these structures as possible
encodings of a richer source, but projection alone does not force strings,
branes, a worldsheet, a critical dimension, modular invariance, or a
graviton. This paper separates the internal results of perturbative string
theory from their proposed interpretation in Modal Triplet Theory. A valid
MTT-to-string bridge must construct target fields, a worldsheet action and
measure, Weyl and modular consistency, anomaly cancellation, factorization,
and the relevant observable dictionary from the same selected source. The
q79 Fu--Yau branch is the present compatible candidate. Its twelve-row
heterotic worldsheet contract has five available rows, two partial rows, and
five open rows. The resulting claim is a structured string-encoding program,
not a derivation of string theory from projection or a completed ultraviolet
theory.
\end{abstract}

\section*{Version 2 Revision Note}
\begin{description}[leftmargin=3.2cm,style=nextline]
\item[Supersedes] Version 1.0, \emph{A Projection-First Reframing of String
Theory}.
\item[Reason] The original paper presented strings, closed-string gravity,
D-branes, holography, and the rejection of most vacua as consequences of
finite admissibility. Those conclusions require independent worldsheet,
boundary, anomaly, and observable structures.
\item[Resolution] String structures are now conditional realization classes.
The paper distinguishes standard string-theoretic results from an MTT
interpretation and states the current q79 worldsheet contract row by row.
\item[Retained result] Projection-first language remains useful for asking
whether several lower descriptions descend from one selected source and
preserve the same physical data.
\item[Remaining boundary] The physical q79 bundle, exact \((0,2)\) theory,
modular/GSO completion, q79 string-field vertices, infrared completion, and
nonperturbative definition remain open.
\end{description}

\tableofcontents

\section{What is being reframed}

String theory is not merely the assertion that fundamental objects are
one-dimensional. Its predictive structure comes from a linked package:
\begin{itemize}[leftmargin=*]
\item a two-dimensional worldsheet field theory;
\item target-space fields and boundary conditions;
\item gauge fixing, ghosts, and a quantum measure;
\item Weyl and modular consistency;
\item GSO projection, factorization, and anomaly cancellation;
\item a spectrum and observable prescription; and
\item compactification and effective-field-theory limits.
\end{itemize}
Extended carriers are one part of this package. A loop drawn in a source
space is not yet a string worldsheet, and a finite projector is not yet a
worldsheet path integral.

The projection-first proposal is interpretive. It asks whether string
variables can be a faithful lower description of selected upper data:
\[
 \Pi_{\rm ws}:\mathcal X_{\rm MTT}
 \longrightarrow
 \mathcal C_{\rm ws}.
\]
The codomain must contain more than embeddings \(X^\mu(\sigma)\). It must
contain all data needed to define the quantum worldsheet theory and its
target-space interpretation.

\section{What standard string theory establishes}

\subsection{Closed strings and gravity}

Within perturbative closed-string theory, quantization of the worldsheet
produces a massless spin-two state. Consistency of its interactions leads to
general covariance and, at low energy, an Einstein-type target-space action
with additional fields and higher-derivative corrections
\cite{GSW,Polchinski1}. This is a strong internal result of the specified
string construction.

It is not a consequence of a generic many-to-one map. A projection can have
no spin-two sector at all. An MTT bridge must therefore show that its selected
source maps to the relevant closed-string Hilbert space, BRST cohomology,
three-point couplings, and low-energy normalization.

\subsection{Open strings and D-branes}

D-branes are boundary conditions for open strings and carry tension,
Ramond--Ramond charge, worldvolume fields, and consistency conditions. Gauge
fields on brane stacks arise from open-string sectors and Chan--Paton data
\cite{Polchinski2}. Calling a boundary an ``admissibility boundary'' can be a
useful analogy, but it does not supply these charges, spectra, or amplitudes.

To realize a D-brane from MTT one must construct:
\[
\begin{aligned}
 &\text{boundary condition}+\text{open-string state space}+\text{charge}\\
 &\hspace{24mm}+\text{worldvolume action}+\text{anomaly conditions}.
\end{aligned}
\]
Without this contract the correspondence is interpretive only.

\subsection{Criticality and modular consistency}

Critical dimensions in perturbative string theories arise from cancellation
of the worldsheet conformal anomaly for a declared matter and ghost system.
They are not dimensional consequences of a circle, a rank flag, or a
projection. Likewise, modular invariance is a property of a torus partition
function and its state sum. A finite arithmetic symmetry may support a
candidate construction, but it does not replace analytic characters, GSO
projection, or factorization.

\section{What projection-first language can claim}

Projection-first language is strongest when used as a typing discipline.
Suppose an upper object \(\mathcal U\) carries connections, actions, states,
and observables. A string encoding is physically meaningful only if maps
\[
 \Pi_{\rm target},\qquad
 \Pi_{\rm ws},\qquad
 \Pi_{\rm obs}
\]
preserve the structures required on their common domain. For example, a
background equation obtained from a target-space action and a beta function
obtained from the worldsheet must agree after conventions and field
redefinitions are fixed. Similar-looking equations are not enough.

This supports a conditional statement:
\begin{quote}
If one selected upper construction admits a complete, anomaly-free
worldsheet realization and a controlled target-space descent, then string
theory can be a lower encoding of that construction on the common domain.
\end{quote}
It does not support the converse claim that finite admissibility forces
strings or that every consistent string vacuum is physically selected.

\section{Why extended carriers are not automatic}

A nontrivial loop or holonomy requires at least a suitable one-dimensional
cycle. This explains why circles and line phases recur in Fourier analysis,
gauge theory, and string compactification. It does not select a dynamical
string. A point, graph, matrix algebra, particle worldline, higher-dimensional
brane, or noncommutative algebra can also carry finite constraints.

Likewise, indecomposability of a constraint system does not imply an
extended carrier. The MTT saturation paper gives finite countermodels:
complete and locally rigid constraint systems can exist on point-valued
carriers. A string-like realization becomes preferred only after additional
requirements such as continuous loop support, reparametrization symmetry,
worldsheet locality, or anomaly cancellation are imposed.

The proper hierarchy is therefore:
\[
 \text{projection}
 \;\not\Rightarrow\;
 \text{extended carrier}
 \;\not\Rightarrow\;
 \text{worldsheet QFT}
 \;\not\Rightarrow\;
 \text{consistent string vacuum}.
\]
Each arrow requires a construction.

\section{Duality, holography, and the landscape}

\subsection{Duality}

A duality requires two theories, an invertible comparison on the declared
domain, and preservation of observables or amplitudes. The shared-circle
Fourier structure can motivate a \(T\)-duality candidate, but exact
string-theoretic \(T\)-duality also acts on momentum, winding, background
fields, and quantum states. \(S\)-duality requires an additional
strong--weak coupling map. Neither follows from the word ``projection.''

\subsection{Holography}

Holographic duality is not simply the survival of information on a boundary.
It needs two fully defined theories and a quantitative dictionary equating
partition functions, states, or correlators. Projection-first language can
classify bulk and boundary variables, but it does not prove AdS/CFT or a
general holographic principle.

\subsection{Vacua and selection}

The string landscape is a set of candidate compactifications subject to
increasingly strong consistency conditions. MTT may add a source-selection
criterion and correlate data ordinarily chosen independently. Yet
mathematical consistency, local stability, cosmological realization, and
empirical adequacy are different filters. No current MTT theorem eliminates
all but one physical string vacuum.

\section{The selected q79 candidate}

The strongest current string-compatible MTT branch is the q79 Fu--Yau
candidate. Its role should be stated carefully:
\begin{itemize}[leftmargin=*]
\item the q79 arithmetic and finite carrier data are exact at their declared
tier;
\item Fu--Yau and related torus-bundle solutions provide established
Hull--Strominger geometry \cite{FuYau,Fino2019};
\item a compatible finite Green--Schwarz/Bianchi representative and
curvature-level visible cancellation are available;
\item the physical nonpullback visible and hidden bundles and their
continuum HYM realization are not yet selected; and
\item no exact q79 \((0,2)\) SCFT or nonperturbative string definition has
been constructed.
\end{itemize}
The auxiliary literal product \(S^1\times{\rm Lens}\times{\rm Nil}\) is not
the q79 compactification. Circle--Lens--Nil language is retained as a
rank/operator filtration; the physical compactification candidate is the
q79 Fu--Yau branch.

\subsection{The twelve-row worldsheet contract}

\small
\begin{longtable}{@{}p{0.07\textwidth}p{0.54\textwidth}p{0.28\textwidth}@{}}
\toprule
Row & Required object & Current status\\
\midrule
\endfirsthead
\toprule
Row & Required object & Current status\\
\midrule
\endhead
W1 & Time-oriented q79 target branch & Available\\
W2 & Fu--Yau Mukai charge and Green--Schwarz/Bianchi sector & Available\\
W3 & Visible curvature-level Green--Schwarz cancellation & Available at
curvature tier\\
W4 & Critical heterotic central-charge balance & Available universally\\
W5 & q79 low-energy GR and quantum-EFT limit & Available at declared parity
tier\\
W6 & Global Deligne gerbe and Freed--Witten condition on the full visible
cycle set & Open; finite representative only\\
W7 & All-orders-in-\(\alpha'\) q79 target background & Open; current
background is first order\\
W8 & Exact q79 heterotic \((0,2)\) SCFT or all beta functions & Partial;
base GLSM, local anomaly, and topological candidate exist, but the physical
bundle and IR SCFT do not\\
W9 & q79 modular-invariant GSO partition function and factorization &
Partial; finite torsion module and seven-seed induction exist, but analytic
characters and GSO remain open\\
W10 & q79-specific string-field vertices and BV master action & Open;
standard framework exists but is not instantiated on q79\\
W11 & Tadpole, vacuum-shift, infrared, and soft-state completion & Open\\
W12 & All-genus convergence or a nonperturbative definition & Open\\
\bottomrule
\end{longtable}
\normalsize

The authoritative count is therefore five available, two partial, and five
open. Standard heterotic string-field BV machinery exists
\cite{Sen2016}; importing that framework is not the same as constructing its
q79 vertices.

\section{Claim ledger}

\begin{center}
\small
\begin{tabularx}{\textwidth}{@{}p{0.31\textwidth}p{0.18\textwidth}X@{}}
\toprule
Claim & Status & Boundary\\
\midrule
Strings are allowed MTT encodings & Supported & Requires a declared bridge
and domain.\\
Projection forces strings & False in general & Counterexamples include
finite point-valued carriers.\\
Closed strings contain a graviton & Standard string result & Assumes the
quantized closed-string construction.\\
MTT derives the closed-string sector & Open & Needs state-space, BRST, and
coupling intertwiners.\\
D-branes are admissibility boundaries & Interpretive analogy & Charges,
boundary CFT, and worldvolume dynamics are additional.\\
q79 selects a compatible string route & Supported & Fu--Yau heterotic
candidate at partial contract tier.\\
q79 defines a complete worldsheet theory & Not established & Current
contract is \(5\) available, \(2\) partial, \(5\) open.\\
MTT uniquely selects a physical vacuum & Open & Needs complete source,
dynamics, and observable tests.\\
\bottomrule
\end{tabularx}
\end{center}

\section{Discussion}

The projection-first interpretation remains valuable after these
restrictions. It changes the research question from ``are strings literally
fundamental?'' to ``which upper data can a string description encode, and
with what loss?'' That question can be tested map by map.

The strongest route is not to redescribe familiar string motifs in new
language. It is to complete the q79 contract and then compare two
independently constructed descendants of the same source: the four-dimensional
effective theory and the worldsheet theory. Agreement of actions, anomalies,
states, and observables would be substantive evidence for a common encoding.
Failure of a row would identify exactly where the analogy stops.

\section{Conclusion}

String theory is a possible and unusually rich projection-compatible
encoding; it is not a consequence of projection alone. The extended
carrier, worldsheet quantum theory, criticality, branes, dualities, and
observables must each be supplied.

MTT has selected exact finite q79 data and a compatible Fu--Yau route, but
the physical string construction remains incomplete. The next decisive
steps are the visible-hidden HYM bundles, the exact \((0,2)\) infrared
theory, modular/GSO completion, and q79-specific string-field vertices.
Until those are built, the honest result is a sharply organized bridge
program rather than a derivation of string theory or a finished ultraviolet
completion.

\begin{thebibliography}{9}

\bibitem{GSW}
M.~B. Green, J.~H. Schwarz, and E. Witten,
\emph{Superstring Theory}, Vols.~1--2,
Cambridge University Press, 1987.

\bibitem{Polchinski1}
J. Polchinski,
\emph{String Theory}, Vol.~1,
Cambridge University Press, 1998.

\bibitem{Polchinski2}
J. Polchinski,
\emph{String Theory}, Vol.~2,
Cambridge University Press, 1998.

\bibitem{FuYau}
J.-X. Fu and S.-T. Yau,
\emph{The Theory of Superstring with Flux on Non-K\"ahler Manifolds and
the Complex Monge--Amp\`ere Equation},
Journal of Differential Geometry \textbf{78} (2008), 369--428.

\bibitem{Fino2019}
A. Fino, G. Grantcharov, and L. Vezzoni,
\emph{Solutions to the Hull--Strominger System with Torus Symmetry},
Communications in Mathematical Physics \textbf{388} (2021), 947--967.

\bibitem{Sen2016}
A. Sen,
\emph{BV Master Action for Heterotic and Type II String Field Theories},
Journal of High Energy Physics \textbf{02} (2016), 087.

\end{thebibliography}

\end{document}
